\subsection*{1.2. 拓扑斯 $Y\times_s S$}

\subsubsection*{1.2.1.}
设 $S$ 为 Hensel 局部概形, 闭点为 $i:s\hookrightarrow S$. 函子
\begin{equation*}
S'\longmapsto S'\times_Ss=i^*S'
\end{equation*}
给出从有限平展 $S$-概形 $S'$ 所成范畴到有限平展 $s$-概形 $s'$（平展 $k(s)$-代数的谱）所成范畴的等价. 逆函子 $\mathrm{sp}^*$ 从 $s$ 的平展景到 $S$ 的平展景, 给出景态射
\begin{equation*}
\mathrm{sp}:S\longrightarrow s,
\end{equation*}
称为\underline{特殊化态射}（specialization morphism）（参见 SGA~4, VIII, 7）. 若 $s'$ 在 $s$ 上平展, 则称 $\mathrm{sp}^*(s')$ 是由 $S$ 经\underline{剩余域扩张}得到的. 对 $S$ 上的层 $F$, 有函子性同构
\begin{equation}\tag{1.2.1.1}
\mathrm{sp}_*F=i^*F.
\end{equation}
对于 $S$ 的任意开子概形 $j:U\hookrightarrow S$, 或 $S$ 的任意点 $j:\eta\hookrightarrow S$, 我们仍用 $\mathrm{sp}$ 表示复合态射 $\mathrm{sp}\circ j$. 对 $U$（分别对 $\eta$）上的层 $F$, 有
\begin{equation}\tag{1.2.1.2}
\mathrm{sp}_*F=i^*j_*F.
\end{equation}

\subsubsection*{1.2.2.}
设 $S$ 为 Hensel 离散赋值环谱, $s,\eta,\bar s,\bar\eta,I$ 如 0.2.5 所示.
\begin{center}
\begin{tikzcd}
s \arrow[r, hook, "i"] & S & \eta \arrow[l, hook, "j"']
\end{tikzcd}
.\end{center}
下面给出 $s$, $\eta$, $S$ 上的层及特殊化态射的 Galois 理论解释.

\begin{enumerate}[label=(\alph*)]
\item $s$（分别 $\eta$）上的层 $F$ 通过 $F\mapsto F_{\bar s}$（分别 $F\mapsto F_{\bar\eta}$）, 等同于带有 $\operatorname{Gal}(\bar s/s)$（分别 $\operatorname{Gal}(\bar\eta/\eta)$）连续作用的集合（SGA~4, VIII, 2.1, 或 1.1.3）.

\item $S$ 上的层 $F$ 定义了 $s$ 和 $\eta$ 上的层 $F_s=i^*F$ 和 $F_\eta=j^*F$, 以及伴随箭头 $F\longrightarrow j_*j^*F$, 它诱导
\begin{equation*}
\varphi:F_s=i^*F\longrightarrow i^*j_*j^*F
=\mathrm{sp}_*F_\eta.
\end{equation*}
已知（SGA~4, IV, 9.5.4）函子
\begin{equation*}
F\longmapsto(F_s,F_\eta,\varphi)
\end{equation*}
给出从 $S$ 上层所成范畴到三元组 $(F_s,F_\eta,\varphi)$ 所成范畴的等价, 其中 $F_s$ 为 $s$ 上的层, $F_\eta$ 为 $\eta$ 上的层, $\varphi:F_s\longrightarrow i^*j_*F_\eta$.

在 (a) 的等价下, 函子 $i^*j_*$ 等同于“取 $I$-不变部分”函子. 因此 $S$ 上的层等同于如下三元组: 一个带有 $\operatorname{Gal}(\bar s/s)$ 作用的集合 $F_{\bar s}$, 一个带有 $\operatorname{Gal}(\bar\eta/\eta)$ 作用的集合 $F_{\bar\eta}$, 以及一个等变态射 $\varphi:F_{\bar s}\longrightarrow F_{\bar\eta}$.

\item 在 (a) 和 (b) 的对应下, 态射 $\mathrm{sp}$, $j$, $i$ 和 $\mathrm{sp}=\mathrm{sp}\circ j$ 表示如下:
\begin{equation*}
\begin{array}{llcl}
\mathrm{sp}:S\longrightarrow s:\;&
 \mathrm{sp}^*(F_{\bar s})&=&(F_{\bar s},F_{\bar s},\mathrm{Id}),\\[2pt]
&\mathrm{sp}_*(F_{\bar s},F_{\bar\eta},\varphi)&=&F_{\bar s},\\[3pt]
j:\eta\longrightarrow S:\;&
 j^*(F_{\bar s},F_{\bar\eta},\varphi)&=&F_{\bar\eta},\\[2pt]
&j_*(F_{\bar\eta})&=&(F_{\bar\eta}^{I},F_{\bar\eta},
 \text{包含}),\\[3pt]
i:s\longrightarrow S:\;&
 i^*(F_{\bar s},F_{\bar\eta},\varphi)&=&F_{\bar s},\\[2pt]
&i_*(F_{\bar s})&=&(F_{\bar s},\{0\},\text{投影}),\\[3pt]
\mathrm{sp}=\mathrm{sp}\circ j:\eta\longrightarrow s:\;&
 \mathrm{sp}^*(F_{\bar s})&=&F_{\bar s},\\[2pt]
&\mathrm{sp}_*(F_{\bar\eta})&=&F_{\bar\eta}^{I}.
\end{array}
\end{equation*}
\end{enumerate}

\subsubsection*{1.2.3.}
设 $S$ 为 1.2.2 中的 Hensel 离散赋值环谱, $Y$ 为 $s$ 上的概形. 我们须使用 $Y$ 和 $S$ 的平展拓扑斯在 $s$ 的平展拓扑斯上的 2-纤维积 $Y\times_sS$. 由 Giraud~[1], 此 2-纤维积存在, 而且容易看出它具有下述 Galois 理论描述. 不喜欢 2-纤维积的读者可将这一描述作为定义.

\subsubsection*{构造 1.2.4.}
\underline{$Y\times_sS$ 上的层 $F$ 等同于如下三元组 $(F_s,F_\eta,\varphi)$:}
\begin{enumerate}[label=(\alph*)]
\item $F_s$ 是 $Y$ 上的层, 或等价地, 是 $\bar Y=Y\otimes_s\bar s$ 上带有 $\operatorname{Gal}(\bar s/s)$ 连续作用且与 $\operatorname{Gal}(\bar s/s)$ 在 $\bar Y$ 上的作用相容的层 $F_{\bar s}$;

\item $F_\eta$ 是 $\bar Y$ 上带有 $\operatorname{Gal}(\bar\eta/\eta)$ 连续作用的层 $F_{\bar\eta}$, 且此作用与 $\operatorname{Gal}(\bar\eta/\eta)$ 在 $\bar Y$ 上（通过 1.1.1）的作用相容;

\item $\varphi$ 是从 $F_{\bar s}$ 到 $F_{\bar\eta}$ 的等变态射.
\end{enumerate}
类似地, \underline{$Y\times_s\eta$ 上的层 $F_\eta$} 是 1.2.4(b) 所述的对象. 有拓扑斯图
\begin{equation}\tag{1.2.4.1}
\begin{tikzcd}[column sep=large, row sep=small]
Y \arrow[dd] & {} & Y \arrow[r, hook, "i"] \arrow[dd]
  & Y\times_sS \arrow[dd]
  & Y\times_s\eta \arrow[l, hook', "j"'] \arrow[dd] & {} \\
{} & {} \arrow[l, "\mathrm{sp}"'] & {} & {} & {} & {} \\
s & {} & s \arrow[r, hook, "i"] & S
  & \eta \arrow[l, hook', "j"'] & {.}
\end{tikzcd}
\end{equation}
拓扑斯 $Y\times_sS$ 是开子拓扑斯 $Y\times_s\eta$ 及其闭补集 $Y$ 的并. 1.2.2(c) 中的公式对 $\mathrm{sp}:Y\times_sS\longrightarrow Y$（也记为 $\mathrm{pr}_1$）, $j:Y\times_s\eta\hookrightarrow Y\times_sS$, $i:Y\hookrightarrow Y\times_sS$, 及 $\mathrm{sp}:Y\times_s\eta\longrightarrow Y$（也记为 $\mathrm{pr}_1$）仍然成立.

\paragraph{1.2.4.2.}
若采用 1.2.4 作为 $Y\times_sS$ 的定义, 就必须通过直接构造说明, 在唯一同构的意义下, $Y\times_sS$ 不依赖于 $k(\eta)$ 的可分闭包 $k(\bar\eta)$ 的选择. 为此, 将 1.2.4(b) 中的 $F_\eta$ 解释为一个函子: 它对 $k(\eta)$ 的每个可分闭包 $k(\bar\eta)$（此闭包定义 $k(s)$ 的一个可分闭包 $k(\bar s)$）, 指定 $Y\otimes_{k(s)}k(\bar s)$ 上的一个层; 1.2.4(b) 中的 $F_\eta$ 就是此函子在 $\{\bar\eta\}$ 上的限制.
